\documentclass[11pt]{article}
\usepackage{../../../tex/lecturestyle}

\title{Sparse Table}
\author{}
\date{}

\begin{document}
\thispagestyle{fancy}

\lecturetitleband{Data Structures}{Sparse Table}{3}

\section{Overview}
A \emph{sparse table} answers range queries such as ``what is the minimum
element of $a[l..r)$?'' in $O(1)$ time after an $O(n\log n)$ preprocessing
step, provided the underlying operation is \emph{idempotent}
($x \oplus x = x$), as is the case for $\min$, $\max$, $\gcd$, bitwise
\texttt{AND}, and bitwise \texttt{OR}. Unlike a segment tree, a sparse table
supports no updates: the array must be fixed before queries begin. In
exchange, queries are simpler and touch only two precomputed cells.

\section{Definitions}
\begin{definitionbox}
Fix an array $a[0..n)$ and an idempotent, associative, commutative operator
$\oplus$. For $k \ge 0$ and $0 \le i \le n - 2^k$, define
\[
  \mathrm{st}[k][i] \;=\; a[i] \oplus a[i+1] \oplus \cdots \oplus a[i+2^k-1],
\]
the $\oplus$-aggregate of the length-$2^k$ window starting at $i$. Row $k=0$
is just $\mathrm{st}[0][i] = a[i]$.
\end{definitionbox}

\section{The Technique}

\subsection{Build}
Each row is built from the previous one by splitting a length-$2^k$ window
into two length-$2^{k-1}$ halves:
\[
  \mathrm{st}[k][i] \;=\; \mathrm{st}[k-1][i] \;\oplus\; \mathrm{st}[k-1][i + 2^{k-1}].
\]
There are $O(\log n)$ rows and $O(n)$ entries per row, so the whole table is
built in $O(n \log n)$ time and space.

\subsection{Query}
To answer a query on $[l, r)$, let $k = \lfloor \log_2(r-l) \rfloor$, so that
$2^k \le r - l < 2^{k+1}$. Cover $[l, r)$ with two windows of length $2^k$:
one starting at $l$, one ending at $r$ (starting at $r - 2^k$). These two
windows may overlap, but because $\oplus$ is idempotent this does not
double-count:
\[
  \mathrm{answer}(l, r) \;=\; \mathrm{st}[k][l] \;\oplus\; \mathrm{st}[k][r - 2^k].
\]

\begin{theorembox}
The query above returns the correct aggregate of $a[l..r)$ in $O(1)$ time,
using $O(1)$ extra arithmetic beyond computing $k = \lfloor \log_2(r-l)\rfloor$
(which itself can be precomputed for all lengths in $O(n)$).

\emph{Why it is correct:} both windows lie inside $[l,r)$ and together cover
it, since the first starts at $l$ and the second ends at $r$ and each has
length $2^k \ge (r-l)/2$. Their union is exactly $[l,r)$, and their
intersection (if any) is covered redundantly; idempotence
($x \oplus x = x$) guarantees that combining two aggregates whose ranges
overlap still yields the aggregate of the union.
\end{theorembox}

\begin{notebox}
Idempotence is essential. For a \emph{non}-idempotent operator such as sum,
this two-window trick over-counts the overlap, and one needs a different
structure (e.g. a Fenwick tree or prefix sums for sum queries, or a segment
tree for point updates).
\end{notebox}

% FIGURE: sparse_query
\begin{figure}[H]
\centering
\begin{tikzpicture}[x=1.0cm,y=0.85cm]
\foreach \i/\v in {0/5,1/2,2/7,3/3,4/9,5/1,6/6,7/4}{
  \node[lcell] at (\i,0) {\v};
  \node[llabel] at (\i,0.75) {\i};
}
\fill[lhi2] (2.5,-1.65) rectangle (4.5,-0.45);
\draw[lectureaccent,thick] (0.5,-1.0) rectangle (4.5,-0.6);
\node[llabel] at (2.5,-0.8) {$\mathrm{st}[2][1]=2$};
\draw[lecturewarn,thick] (2.5,-1.6) rectangle (6.5,-1.2);
\node[llabel] at (4.5,-1.4) {$\mathrm{st}[2][3]=1$};
\draw[decorate,decoration={brace,amplitude=4pt}] (0.5,1.15) -- (6.5,1.15)
  node[midway,above=4pt,llabel] {query $[1,7)$, \; $k=\lfloor\log_2 6\rfloor=2$};
\node[llabel,anchor=west] at (7.0,-1.0) {$\min(2,1)=1$};
\end{tikzpicture}
\caption{A query on $a=(5,2,7,3,9,1,6,4)$ with $\oplus=\min$. The range has length 6, so $k=2$; the two length-4 windows start at $l=1$ and at $r-2^k=3$ and overlap on positions 3 and 4 (shaded), which is harmless because $\min$ is idempotent.}\label{fig:sparse_query}
\end{figure}

\section{Complexity}
\begin{itemize}[leftmargin=1.5em]
  \item Build: $O(n \log n)$ time, $O(n \log n)$ space.
  \item Query: $O(1)$ time (after $O(n)$ precomputation of $\lfloor\log_2 x\rfloor$).
  \item No updates supported.
\end{itemize}

\section{Reference Implementations}

\codenote
\subsection{Rust}
\corelisting{Rust}{example.rs}

\subsection{Python}
\corelisting{Python}{example.py}

\section{Worked Example}
\begin{examplebox}
Let $a = [5, 2, 4, 7, 1, 3, 6]$ with $\oplus = \min$. To answer a range-min
query on $[1, 6)$ (values $2,4,7,1,3$, whose minimum is $1$): the length is
$5$, so $k = \lfloor \log_2 5 \rfloor = 2$, i.e. windows of length $4$. The
first window starts at $1$: $\mathrm{st}[2][1] = \min(2,4,7,1) = 1$. The
second window ends at $6$, so it starts at $6-4=2$:
$\mathrm{st}[2][2] = \min(4,7,1,3) = 1$. Combining, $1 \oplus 1 = 1$, the
correct answer.
\end{examplebox}

\section{Further Reading}
\begin{itemize}[leftmargin=1.5em]
  \item AtCoder Algorithm Lectures index: \url{https://info.atcoder.jp/entry/algorithm_lectures/index}
\end{itemize}

\end{document}
